\newcommand{\ModuleID}{EL-II.12}
\newcommand{\ModuleTitle}{Gerund, Gerundive, Supine, and Passive Periphrastic}
\newcommand{\ModuleStage}{Stage II — Core Grammar}
\newcommand{\ModuleUnit}{II.12}
\newcommand{\ModuleOutcome}{Distinguish verbal nouns and adjectives, convert object-taking gerunds to gerundives, and analyze purpose and obligation constructions}
\newcommand{\ModulePrerequisites}{EL-II.11 mastery: control commands, prohibitions, wishes, and fear-clause polarity}
\newcommand{\ModuleSessionOne}{Build the gerund-to-gerundive and purpose-form ladder, then complete Worksheet II.45 as the guided application.}
\newcommand{\ModuleSessionTwo}{Distinguish future participle, gerundive, and supines in context; map personal and impersonal obligation with dative agent; then complete Worksheet II.48 independently.}
\newcommand{\ModulePrintedPractice}{Worksheets II.45 (guided Variant A) and II.48 (independent Variant D), with terminal keys}
\newcommand{\ModuleExtensionPractice}{Worksheets II.46 and II.47 remain optional remediation or delayed-recall variants}
\newcommand{\ModuleNext}{EL-II.13, Place, Time, Means, Manner, and Cause}
\newcommand{\ModuleMastery}{Identify at least eleven of twelve mixed gerunds, gerundives, future participles, and supines; convert four object-taking gerunds with exact agreement; and compose personal and impersonal passive periphrastics with correct dative agent.}
\newcommand{\ModuleDelayNote}{}
\newcommand{\ModuleReferenceFocus}{%
  \begin{itemize}
    \item gerund as neuter singular verbal noun in its surviving oblique cases;
    \item gerundive as a declining verbal adjective agreeing with a noun;
    \item preferred gerund-to-gerundive conversion when a direct object is present;
    \item accusative and ablative supines; and
    \item passive periphrastic obligation with personal or impersonal construction and dative agent.
  \end{itemize}}
\newcommand{\ModuleContrast}{A gerund is a verbal noun; a gerundive agrees as a verbal adjective. A future active participle expresses prospective active action, while a gerundive can contribute passive necessity; similar-looking stems do not make them interchangeable.}
\newcommand{\ModuleLessonSource}{\CourseRoot/shared/content/volume2/all}
\newcommand{\ModuleExerciseSource}{\CourseRoot/shared/exercises/volume2/all}
\newcommand{\ModuleWorkedBank}{II}
\newcommand{\ModuleExerciseBank}{II}
\newcommand{\ModuleWorksheetVariants}{A,D}
\newcommand{\ModuleScopeNote}{This packet teaches the surviving forms and central prose uses required by the curriculum; lexical restrictions on individual supines remain source-specific.}
\newcommand{\ModuleEnrichment}{%
  \SubmoduleExtension
    {Verbal nouns and adjectives}
    {A gerund-to-gerundive ladder separates the unchanged proposition from the new agreement structure.}
    {Write the lesson's object-taking gerund phrase on the first rung. On the second, move the object into the gerundive's agreeing noun phrase; label case, number, gender, and unchanged verbal relation. Add a third rung showing a passive-periphrastic obligation.}
    {Define gerund and gerundive by part of speech and agreement behavior.}
    {The gerund is a neuter singular verbal noun; the gerundive is a declining verbal adjective. The conversion must preserve the governing preposition or case function and the proposition while making the gerundive agree exactly with its noun.}
  \SubmoduleExtension
    {Reading the forms in context}
    {Form, attachment, and governor distinguish gerund, gerundive, future participle, and supine.}
    {For the lesson's mixed examples, fill four fields: morphological form, agreeing controller if any, governing construction, and active/passive or obligation contribution.}
    {State the decisive difference between a future active participle and a gerundive.}
    {A future active participle agrees with its controller and presents prospective active action; a gerundive is a passive verbal adjective and can participate in obligation. Supines and gerunds occupy their limited noun-like cases and do not agree with a controller.}
  \SubmoduleExtension
    {Writing laboratory}
    {An obligation ladder makes subject and agent visible in personal and impersonal periphrastics.}
    {Complete the lesson's personal and impersonal obligation prompts. Draw an arrow from gerundive to grammatical subject and another from the dative agent to the responsible person; mark where no personal subject is expressed.}
    {Explain why the dative agent is not the grammatical subject of the passive periphrastic.}
    {The gerundive agrees with the grammatical subject in the personal construction; the dative names the person upon whom the obligation rests. In the impersonal construction the verbal adjective is neuter singular and no personal nominative subject is supplied.}
}
